\documentclass[11pt]{article}

\usepackage[utf8]{inputenc}
\usepackage[T1]{fontenc}
\usepackage{lmodern}
\usepackage{geometry}
\usepackage{amsmath,amssymb,amsfonts,amsthm,mathtools}
\usepackage{bm}
\usepackage{enumitem}
\usepackage{microtype}
\usepackage{hyperref}
\usepackage{titlesec}
\usepackage{booktabs}
\usepackage{array}
\usepackage{setspace}
\usepackage{mathrsfs}
\usepackage{physics}

\geometry{margin=1in}
\hypersetup{colorlinks=true, linkcolor=black, urlcolor=blue, citecolor=black}
\setlist[enumerate]{topsep=0.4em,itemsep=0.35em}
\setlist[itemize]{topsep=0.3em,itemsep=0.25em}
\titleformat{\section}{\large\bfseries}{\thesection.}{0.5em}{}
\titleformat{\subsection}{\normalsize\bfseries}{\thesubsection.}{0.5em}{}
\setstretch{1.06}

\newcommand{\Aether}{\AE{}ther}
\newcommand{\Aflow}{\AE{}ther-flow}
\newcommand{\TheoryName}{\Aflow{} Interpretation of Relativity}
\newcommand{\TheoryProgram}{\Aether{} / \Aflow{} framework}
\newcommand{\ObsBurden}{\mathcal{B}_{\mathrm{obs}}}
\newcommand{\ObserverClass}[1]{\left[#1\right]_{\mathrm{obs}}}

\newtheorem{definition}{Definition}
\newtheorem{proposition}{Proposition}
\newtheorem{remark}{Remark}
\newtheorem{corollary}{Corollary}
\newtheorem{theorem}{Theorem}

\title{\textbf{The \TheoryName{}}\\[0.4em]
\large Post-Bridge Derivation-Gate Status Note}
\author{}
\date{}

\begin{document}

\maketitle

\begin{abstract}
The benchmark derivation-gates note fixed the five observer-level gates that any honest substrate derivation of GR must satisfy. The continuity-Hessian bridge theorem-building note now changes the routing status of those gates on the active one-metric charge-polarization line. On the branch-locked active completed branch, the direct factorization theorem proves
\[
\ObserverClass{\ObsBurden}=0,
\]
so the no-genuine-observer-side-remainder gate is no longer the active bottleneck on that same-package observer-side line.

This note records the benchmark-facing consequence. Gates (1)--(4) remain the fixed benchmark conditions of the project and are preserved on the active one-metric package, while Gate (5) is now settled on the active completed branch at same-package observer-side scope. The honest remaining burden is therefore not another package-preserving same-output note. It is the deeper derivational burden of explaining why the benchmark observer package itself arises from explicit substrate structure in the sense relevant to \TheoryName{}. An explicit same-package rewritten-metric reformulation is secondary and matters only if a second exact re-expression is specifically wanted.
\end{abstract}

\section{Introduction}

The active manuscript sequence of \TheoryName{} already has a fixed public front door: the overview, the exact-closure note, and the five core companion manuscripts on foundations, dynamics, consistency, relativistic recovery, and geometry. That package fixes the benchmark theory claim. It is not the place where a substrate derivation is claimed complete.

The derivation-gates note then made the benchmark target explicit: any candidate derivation must recover Einsteinian observer dynamics through one operative metric with universal matter coupling, preserve the ordinary GR gauge and two-tensor content, reproduce exact null/proper-time/redshift recovery, and leave no genuine observer-side remainder. At the time that note was written, the live bottleneck on the active bridge line was still the observer-side remainder gate.

The continuity-Hessian bridge theorem-building note changes that routing status. On the sharpened continuity-Hessian package, the direct bridge is now completed on the branch-locked active completed branch, and the observer-side burden class closes there. A short benchmark-facing status note is therefore needed so that the project does not keep speaking as if Gate (5) were still the active bottleneck on this same line.

\paragraph{Framework and claim boundary.}
\TheoryProgram{} denotes the broader \Aether{} / \Aflow{} framework built on the claims that the \Aether{} is the underlying four-dimensional substrate of reality, the \Aflow{} is its intrinsic ordered motion, observed three-dimensional space is the local experiential slice of that deeper substrate, S-time is the experienced order of change arising from matter, light, and the \Aflow{}, and observed expansion is the three-dimensional appearance of deeper four-dimensional ordered motion. Within that framework, \TheoryName{} denotes the exact-closure theory in which the effective gravitational dynamics are adopted to be exactly Einsteinian with universal matter coupling. In the active sequence, \emph{adoption} means use of that established relativistic dynamics without claiming substrate derivation, while \emph{derivation} is reserved for a first-principles recovery from explicit substrate variables.

The present note stays strictly inside that boundary. It does not claim that GR has already been derived from substrate microphysics. Its narrower task is benchmark facing: state plainly what the new continuity-Hessian theorem settles, record the resulting gate status, and identify the next honest open burden.

\section{What the Bridge Theorem Changes}

\begin{definition}[Post-bridge fixed inputs]
The post-bridge status recorded here uses three fixed inputs:
\begin{enumerate}
    \item the benchmark exact-closure package, which fixes the public relativistic target of the repository;
    \item the benchmark derivation-gates note, which states the five non-negotiable gates for an honest substrate derivation;
    \item the continuity-Hessian bridge theorem-building note, which proves the direct factorization theorem and the branch-locked vanishing
    \[
    \ObserverClass{\ObsBurden}=0
    \]
    on the active completed branch of the sharpened one-metric package.
\end{enumerate}
\end{definition}

\begin{proposition}[Gate (5) is no longer the active bottleneck on the same package]
On the active completed branch of the sharpened continuity-Hessian package, the no-genuine-observer-side-remainder gate is settled at same-package observer-side scope.
\end{proposition}

\begin{proof}
The completed observer equation on the active one-metric line already isolates the residual observer-side burden into the single class \(\ObsBurden\), modulo the previously separated gauge-bookkeeping, gauge-exact, and constraint-exact sectors. The continuity-Hessian bridge theorem-building note then proves the direct factorization theorem and, under branch locking on the active completed branch, the exact vanishing
\[
\ObserverClass{\ObsBurden}=0.
\]
That is precisely the observer-level content required for the benchmark no-genuine-remainder gate on this same package. Therefore Gate (5) is no longer the active bottleneck on that branch.
\end{proof}

\begin{remark}
This proposition is narrower than a completed first-principles derivation. It says that the current same-package observer-side remainder burden is closed on the active completed branch. It does not yet explain from deeper substrate structure why that observer package itself should arise.
\end{remark}

\section{Benchmark Gate Status After the Bridge Theorem}

\begin{definition}[Status language]
The gate statuses in Table \ref{tab:gates} use the following phrases.
\begin{enumerate}
    \item \textbf{Settled on the active completed branch:} the current one-metric package now proves the required observer-side statement on the active completed branch.
    \item \textbf{Fixed and preserved on the active package:} the relevant benchmark feature is already built into the current one-metric package and remains intact on the active line, but the deeper question of why that package arises from substrate structure is still open.
    \item \textbf{Open as deeper derivational burden:} what remains is not another same-package observer-side rewrite, but an upstream substrate-level account strong enough to make the benchmark package itself answerable as derivation rather than adoption plus preservation.
\end{enumerate}
\end{definition}

\begin{table}[ht]
\centering
\renewcommand{\arraystretch}{1.2}
\begin{tabular}{>{\raggedright\arraybackslash}p{0.22\textwidth} >{\raggedright\arraybackslash}p{0.34\textwidth} >{\raggedright\arraybackslash}p{0.34\textwidth}}
\toprule
\textbf{Gate} & \textbf{Current status on the active one-metric line} & \textbf{Honest remaining burden} \\
\midrule
Einsteinian observer dynamics & Settled on the active completed branch at current observer-side package scope: once the burden class closes, the completed observer equation reduces to the Einsteinian observer law already isolated on that same package. & Explain from deeper substrate structure why that Einsteinian observer package arises, rather than merely preserving an already recorded package. \\
One operative metric with universal matter coupling & Fixed and preserved on the active one-metric package. & Derive that same one-metric universal-coupling structure from explicit substrate variables rather than taking it as preserved package data. \\
Ordinary GR gauge and two-tensor content & Fixed and preserved on the active package at benchmark scope. & Supply the deeper derivational account that makes that same low-energy gauge and two-tensor content emerge from substrate structure. \\
Exact null, proper-time, and redshift recovery & Fixed and preserved on the active package through the same operative metric. & Explain from deeper substrate structure why that same operative metric carries the full benchmark relativistic recovery package. \\
No genuine observer-side remainder & Settled on the active completed branch by the direct continuity-Hessian closure theorem, which proves \(\ObserverClass{\ObsBurden}=0\). & No further same-package observer-side remainder note is primary by default; the remaining burden shifts upstream to actual substrate derivation of the benchmark package itself. \\
\bottomrule
\end{tabular}
\caption{Post-bridge benchmark derivation-gate status on the active one-metric charge-polarization line.}
\label{tab:gates}
\end{table}

\begin{theorem}[Primary burden after post-bridge gate status]
After the continuity-Hessian direct theorem and the gate-status restatement of Table \ref{tab:gates}, the next primary benchmark-facing burden is no longer another package-preserving same-output note on the same completed observer package. It is the deeper derivational burden of recovering that benchmark package itself from explicit substrate structure, in a form answerable to all five benchmark gates as derivation rather than adoption plus preservation.
\end{theorem}

\begin{proof}
By Proposition 1 and the Gate (5) row of Table \ref{tab:gates}, the no-genuine-observer-side-remainder gate is no longer the active bottleneck on the branch-locked active completed branch. The remaining rows show that the other benchmark features are already fixed and preserved on the active one-metric package. Therefore another note that merely preserves the same package and reproduces the same observer output does not change the benchmark-facing status. What remains is the upstream derivational burden of explaining why that package itself arises from deeper substrate structure in the sense required by the derivation gates. That is exactly the distinction between adoption and derivation already fixed by the project claim boundary.
\end{proof}

\begin{corollary}[Status of the rewritten-metric route]
An explicit same-package rewritten-metric reformulation is no longer the default next benchmark-facing necessity. It matters only as an optional second exact re-expression of the same closure result.
\end{corollary}

\begin{proof}
The direct same-package closure result is already in hand on the active completed branch. Therefore an explicit rewritten metric does not remove a still-open direct observer-side burden on that branch. It can matter only if one wants a second exact formulation of the same observer-side closure statement.
\end{proof}

\section{Conclusion}

The positive result recorded here is benchmark facing and narrow. The active one-metric charge-polarization line no longer needs to be described as if Gate (5) were still its unresolved front-line bottleneck. On the branch-locked continuity-Hessian package, the direct closure theorem now settles the no-genuine-observer-side-remainder gate on the active completed branch, and therefore also removes the need for another default same-package direct-burden note.

The claim boundary remains disciplined. This is not yet a completed first-principles substrate derivation of GR. The benchmark package still stands as exact relativistic adoption, and the present bridge result is a same-package observer-side closure theorem on top of that package. The open burden has moved upstream: the project now needs a deeper derivational manuscript that explains why the benchmark one-metric observer package itself arises from explicit substrate structure in a way answerable to the five benchmark gates.

That is the correct routing consequence of the new theorem. A merely package-preserving same-output deeper-origin relay is no longer the default primary move. An explicit same-package rewritten-metric reformulation may still be worth recording later, but only as an optional second exact expression of a closure result that is already in hand.

\end{document}
